\begin{abstract}
Construímos uma biblioteca enumerável, sem repetições, de segmentos afins de inclinação racional não nula ancorados em células diádicas. Cada endereço recupera exatamente a célula e a função afim. Montagens contínuas finitas são densas em $C(K)$ e admitem representação exata por redes ReLU rasas racionais com quebras diádicas. Para funções $C^2$, a racionalização preserva a constante $1/8$ do erro interpolatório uniforme. O refinamento adaptativo combina majorantes intervalares racionais da segunda derivada e racionalização nodal, com correção, término e tamanho controlados sob as hipóteses explicitadas. A contagem por altura, os custos de codificação, a entropia e a seleção em catálogos finitos são desenvolvidos no segundo artigo da série.

\smallskip\noindent\textbf{Palavras-chave:} aproximação linear por partes; células diádicas; altura aritmética; endereçamento; refinamento certificado; redes ReLU.
\end{abstract}

\begingroup\small
\begin{center}\textbf{Abstract}\end{center}
\noindent We construct a repetition-free enumerable library of affine segments with nonzero rational slopes anchored on dyadic cells. Each address recovers the cell and the affine function exactly. Finite continuous assemblies are dense in $C(K)$ and admit exact representations as rational shallow ReLU networks with dyadic breakpoints. For $C^2$ functions, rationalization preserves the constant $1/8$ in uniform interpolation error. Adaptive refinement combines rational interval bounds for the second derivative and nodal rationalization, with correctness, termination and size bounds under the stated hypotheses. Height counting, coding costs, entropy and selection in finite catalogs are developed in the second article of the series.

\smallskip\noindent\textbf{Keywords:} piecewise-linear approximation; dyadic cells; arithmetic height; addressing; certified refinement; ReLU networks.
\endgroup
